\documentclass[10pt,a4paper]{amsart}
\usepackage[margin=19mm]{geometry}
\usepackage{amsmath,amssymb,mathtools}
\usepackage[scr=rsfs,cal=euler]{mathalfa}
\usepackage{xfrac}
\usepackage[unicode,hidelinks,bookmarks=false]{hyperref}
\newcommand{\ie}{i.e.\ }
\newcommand{\cf}{cf.\ }
\newcommand{\resp}{respectively\ }
\newcommand{\Subsection}{\subsection}
\providecommand{\nobreakdash}{\nobreak\hbox{-}}
\setlength{\emergencystretch}{2em}
\multlinegap0pt
\newcommand{\K}{\mathbb{K}}
\usepackage[shortlabels]{enumitem}
\usepackage{amssymb}
\newtheorem{theorem}{Theorem}
\newtheorem{corollary}{Corollary}
\title{\bf On a Lemma of Singh and Kumar}
\author{Rishu Garg, Jitender Singh}
\date{Source: arXiv:2511.21427v2 (CC BY 4.0)}
\begin{document}
\maketitle

\noindent\emph{Source and licence.} \bf On a Lemma of Singh and Kumar. Version arXiv:2511.21427v2 (CC BY 4.0), distributed by arXiv under the Creative Commons Attribution 4.0 International licence, http://creativecommons.org/licenses/by/4.0/, as recorded in the arXiv licence metadata for this exact version and shown on the abstract page https://arxiv.org/abs/2511.21427v2. Full licence notice: \texttt{licenses/arxiv-2511.21427-CC-BY-4.0.txt}.\par\smallskip

\noindent\textit{Editorial scope.} Counted unit: the article's corollary c:1 (source0.tex lines 132--141). The article's complete original proof of it is delivered with it (source0.tex lines 240--242, one \texttt{proof} environment of 3 lines, which the article places away from the statement). No mathematics is written, repaired or re-derived: the statement and the proof are the article's own lines, in the article's own notation, and no retained line is substituted or re-flowed.  Retained uncounted as the source-local context the proof needs: lines 121--130.  Nothing was omitted from any retained span, so the package carries no omission record.  No retained line contains a citation, so the wrapper carries no bibliography and no bibliography entry.  No retained line contains a figure environment, an image-inclusion command or drawing code, so no image is delivered and the package declares no asset dependency.  Editorial formatting only, in addition: an AMS \texttt{amsart} preamble that restates the article's own declarations and macros used by the retained lines, namely the environments \texttt{theorem}, \texttt{corollary} on separate counters, together with the standard packages those lines need.  The retained environments print in this document as Theorem 1, Corollary 1 in order of appearance, whereas the article prints the same items with the numbering of its own full document.  The complete native source of the article as distributed in the arXiv e-print for this exact version is bundled unchanged as \texttt{sources/189864/source0.tex}.
\begin{theorem}\label{th:1}
Let $v$ be a Krull valuation of arbitrary rank of a field $\K$ with the value group $G_v$. Let $f=a_0+a_1 z+\cdots+a_n z^n \in \K[z]$ be a polynomial of degree $n$. Suppose there exist smallest indices $j$ and $k$ with $1\leq k+1\leq j\leq n$ for which the following conditions are satisfied.
\begin{enumerate}[label=$(\roman*)$]
		\item[$(i)$] $v(a_j)=0$,
\item[$(ii)$] $\frac{v(a_k)}{j-k}<\frac{v(a_i)}{j-i}$ for all indices $i$ with $0\leq i\leq j-1$ and $i\neq k$,
\item[(iii)] If $j<n$, then $\frac{v(a_k)}{j-k}>\frac{v(a_i)}{j-i}$ for all indices $i$ with $j+1\leq i\leq n$,
\item[$(iv)$] $v(a_k)\notin dG_v$ for any positive divisor $d>1$ of $j-k$.
	\end{enumerate}
Then any factorization $f(z)=f_1(z)f_2(z)$ of $f$ in $\K[z]$ has a factor of degree $\leq n-j+k$. In particular, if $j=n$ and $k=0$, then the polynomial $f$ is irreducible in $\K[z]$.
\end{theorem}

\begin{corollary}\label{c:1}
Let $v$ be a valuation of a field $\K$ with its value group $\mathbb{Z}$. Let $f=a_0+a_1z+\cdots +a_nz^n \in \K[z]$ be a polynomial of degree $n$. Suppose there exist smallest indices $j$ and $k$ with $1\leq k+1\leq j\leq n$ for which the following conditions are satisfied.
	\begin{enumerate}[label=$(\roman*)$]
		\item $v(a_j)=0$,
		\item $\frac{v(a_k)}{j-k}<\frac{v(a_i)}{j-i}$ for all indices $i$ with $0\leq i\leq j-1$ and $i\neq k$,
		\item If $j<n$, then $\frac{v(a_k)}{j-k}>\frac{v(a_i)}{j-i}$ for all indices $i$ with $j+1\leq i\leq n$,
		\item $\gcd(v(a_k),j-k)=1$.
	\end{enumerate}	
	Then any factorization $f(z)=f_1(z)f_2(z)$ of $f$ in $\K[z]$ has a factor of degree $\leq n-j+k$. In particular, if $j=n$ and $k=0$, then the polynomial $f$ is irreducible in $\K[z]$.
\end{corollary}

\begin{proof}[\bf Proof of Corollary \ref{c:1}]
Observe that the hypothesis $(iv)$ of the corollary implies that for any divisor $d>1$ of $j-k$, $v(a_k)\not\in dG_v=d\mathbb{Z}$, since otherwise, $d$ divides $v(a_k)$, and so, $d$ divides $\gcd(v(a_k),j-k)=1$, which is a contradiction. So, the polynomial $f$ satisfies the hypotheses of Theorem \ref{th:1} for $G_v=\mathbb{Z}$, and the corollary follows.
\end{proof}

\end{document}
